% Distortion (mosfet-distortion): the set-up a person would wire, semantically the deck's.
%   VDD vdd 0 20, RL vdd d 1000       copy A: the supply and the drain load                     -> source V_DD (set) + resistor R_L
%   RS s sx 100, CS sx 0 1000u        the source resistor to a node held at AC ground            -> resistor R_S + capacitor C_S
%   IS sx vss 0.01, VSS vss 0 -10     the sink to the negative rail                              -> current source I_S (set) + source V_SS (set)
%   RIN in g 1000, VIN in 0 SIN(…)    the 1 kHz sine through 1 kΩ                                -> resistor R_in + source v_in (set)
%   M1 d g s DUT, V(d) read           the common-source stage; its output                       -> nmos, voltmeter v_out (measured)
%   M1 vdd g s, IS s vss 0.1, CL s l 1000u, RL l 0 100   copy B, the follower deck: the drain on the rail, the load AC-coupled   -> voltmeter v_out on the load
%   The follower is a second deck; the gain-finding probe run is not drawn.
\documentclass[tikz, border=10pt]{standalone}
\usepackage{bench}
\begin{document}
\begin{tikzpicture}
\benchtitle{0,7.8}{Distortion}
\def\gnd{2.0}
% the transistor
\node[nmos] (QA) at (4,3) {};
\node[anchor=west, font=\small, align=left] at ($(QA.D)+(0.15,-0.55)$) {A};
% copy A: the drain load up to the rail
\draw (QA.D) -- ++(0,0.4) coordinate (Adj) to[R, l_={$R_L$}] ++(0,1.6) coordinate (Atop);
\coordinate (Avtop) at ($(Atop)+(4.2,0)$);
\draw (Atop) -- (Avtop);
\draw[set] (Avtop) to[V, l={$V_{DD}$}] (Avtop |- 0,\gnd) node[ground, black]{};
\nodename{anchor=south}{$(Atop)+(1.0,0.05)$}{vdd}
\nodename{anchor=east}{$(QA.D)+(-0.15,0.2)$}{d}
% the gate: the source resistance and the signal source
\coordinate (Ag) at (QA.G);
\draw (Ag) -- ++(-0.3,0) coordinate (Agj) to[R, l_={$R_{in}$}] ++(-1.8,0) coordinate (Avin);
\draw[set] (Avin) to[V, l_=$v_{in}$] (Avin |- 0,\gnd) node[ground, black]{};
\nodename{anchor=south}{$(Ag)+(-0.15,0.05)$}{g}
\nodename{anchor=south east}{$(Avin)+(-0.05,0.05)$}{in}
% the source: the resistor and the sink to the negative rail
\draw (QA.S) to[R, l={$R_S$}] ++(0,-1.2) coordinate (Ax);
\draw[set] (Ax) to[I, l={$I_S$}] ++(0,-1.3) coordinate (Avee);
\draw[set] (Avee) to[V, l=$V_{SS}$] ++(0,-1.3) coordinate (Abot) node[ground, black]{};
\nodename{anchor=south west}{$(Avee)+(0.1,0.05)$}{vss}
\nodename{anchor=west}{$(Abot)+(0.3,0.1)$}{0}
\nodename{anchor=east}{$(QA.S)+(-0.1,-0.2)$}{s}
% the bypass
\draw (Ax) -- ++(-1.3,0) coordinate (Acx) to[C, l_={$C_S$}] (Acx |- 0,\gnd) node[ground]{};
% the voltmeter on Adj
\coordinate (Avo) at ($(Adj)+(2.6,0)$);
\draw (Adj) -- (Avo);
\voltmeter{Avo |- 0,\gnd}{Avo}{$v_{out}$}
\draw (Avo |- 0,\gnd) node[ground]{};
\nodename{anchor=south west}{$(Ax)+(0.1,0.05)$}{sx}
% the transistor
\node[nmos] (QB) at (14.5,3) {};
\node[anchor=west, font=\small, align=left] at ($(QB.D)+(0.15,-0.55)$) {B};
% copy B: the drain on the rail
\draw (QB.D) -- ++(0,2.0) coordinate (Btop);
\coordinate (Bvtop) at ($(Btop)+(5.6,0)$);
\draw (Btop) -- (Bvtop);
\draw[set] (Bvtop) to[V, l={$V_{DD}$}] (Bvtop |- 0,\gnd) node[ground, black]{};
\nodename{anchor=south}{$(Btop)+(1.0,0.05)$}{vdd}
% the gate: the source resistance and the signal source
\coordinate (Bg) at (QB.G);
\draw (Bg) -- ++(-0.3,0) coordinate (Bgj) to[R, l_={$R_{in}$}] ++(-1.8,0) coordinate (Bvin);
\draw[set] (Bvin) to[V, l_=$v_{in}$] (Bvin |- 0,\gnd) node[ground, black]{};
\nodename{anchor=south}{$(Bg)+(-0.15,0.05)$}{g}
\nodename{anchor=south east}{$(Bvin)+(-0.05,0.05)$}{in}
% the source: the resistor and the sink to the negative rail
\draw (QB.S) -- ++(0,-0.5) coordinate (Bx);
\draw[set] (Bx) to[I, l={$I_S$}] ++(0,-1.3) coordinate (Bvee);
\draw[set] (Bvee) to[V, l=$V_{SS}$] ++(0,-1.3) coordinate (Bbot) node[ground, black]{};
\nodename{anchor=south west}{$(Bvee)+(0.1,0.05)$}{vss}
\nodename{anchor=west}{$(Bbot)+(0.3,0.1)$}{0}
\nodename{anchor=west}{$(QB.S)+(0.1,-0.3)$}{s}
% the load, AC-coupled
\coordinate (Bsj) at ($(QB.S)+(0,-0.5)$);
\draw (Bsj) -- ++(1.0,0) to[C, l={$C_L$}] ++(1.6,0) coordinate (Blj);
\draw (Blj) to[R, l_={$R_L$}] (Blj |- 0,\gnd) node[ground]{};
% the voltmeter on Blj
\coordinate (Bvo) at ($(Blj)+(1.8,0)$);
\draw (Blj) -- (Bvo);
\voltmeter{Bvo |- 0,\gnd}{Bvo}{$v_{out}$}
\draw (Bvo |- 0,\gnd) node[ground]{};
\nodename{anchor=south east}{$(Blj)+(-0.1,0.05)$}{l}
% the definitions
\definitions{0,-2.4}{%
  A: $V_{DD} = 20$ V;\quad $I_S = 10$ mA to $V_{SS} = -10$ V;\quad $R_S = 100$ \ohm, bypassed by $C_S = 1000$ \textmu F;\quad $R_L = 1$ k\ohm;\quad $R_{in} = 1$ k\ohm;\quad $f = 1$ kHz\\[2pt]
  B: $V_{DD} = 25$ V;\quad $I_S = 100$ mA;\quad $R_L = 100$ \ohm{} through $C_L = 1000$ \textmu F;\quad $R_{in} = 100$ \ohm{} at 1 and 20 kHz, and 1 k\ohm{} at 20 kHz\\[2pt]
  $v_{in}$ = sine, its level set by a 20 mV probe run for 1 V peak at the output\\[2pt]
  read $v_{out}$ over the last four of ten cycles: THD, HD2, HD3 by Goertzel%ENDDEFS
}
\end{tikzpicture}
\end{document}
